\begin{tcolorbox}
\textbf{هدف این مرحله}

یادگیری SQLite و ذخیره داده سنسور.

\textbf{خروجی:}
Python $\rightarrow$ SQLite $\rightarrow$ SELECT
\end{tcolorbox}

% =========================================================
% عنوان پازل
% =========================================================

\section{هدف این بخش}

در این بخش هدف ما یادگیری و پیاده‌سازی ............ است.

هدف این نیست که تمام مباحث مربوط به ............ را بررسی کنیم،
بلکه فقط مفاهیم مورد نیاز پروژه را یاد می‌گیریم.


% =========================================================

\section{این پازل چیست؟}

توضیح کوتاه و کاربردی درباره موضوع.

در این قسمت فقط به اندازه‌ای توضیح می‌دهیم که بدانیم
این فناوری چه کاری انجام می‌دهد و چرا در پروژه به آن نیاز داریم.


% =========================================================

\section{چرا در این پروژه به آن نیاز داریم؟}

توضیح ارتباط این فناوری با پروژه اصلی.


% =========================================================

\section{پیش‌نیازها}

برای انجام این بخش باید موارد زیر را بلد باشیم:

\begin{itemize}

    \item مورد اول
    \item مورد دوم
    \item مورد سوم

\end{itemize}


% =========================================================

\section{مباحثی که باید یاد بگیریم}

\subsection{مبحث اول}

توضیح کوتاه.

\subsection{مبحث دوم}

توضیح کوتاه.

\subsection{مبحث سوم}

توضیح کوتاه.


% =========================================================

\section{منابع مطالعه}

در این قسمت منابع مفید برای مطالعه بیشتر قرار می‌گیرند.

\begin{itemize}

    \item مستندات رسمی
    \item آموزش پیشنهادی
    \item مقاله یا کتاب در صورت نیاز

\end{itemize}


% =========================================================

\section{پیاده‌سازی اولیه}

در این قسمت یک مثال کاملاً مستقل و کوچک پیاده‌سازی می‌کنیم.

هدف این مثال، یادگیری فناوری است و هنوز الزاماً
بخشی از پروژه نهایی نیست.


% =========================================================

\section{آزمایش و تست}

در این قسمت مشخص می‌کنیم چگونه بفهمیم
پیاده‌سازی ما درست کار می‌کند.


\subsection{آزمایش اول}

شرح آزمایش.


\subsection{نتیجه مورد انتظار}

نتیجه‌ای که باید مشاهده شود.


% =========================================================

\section{اتصال به پروژه اصلی}

در این قسمت مشخص می‌کنیم این پازل
چگونه وارد معماری اصلی پروژه می‌شود.


% =========================================================

\section{خروجی این مرحله}

در پایان این بخش باید بتوانیم:

\begin{itemize}

    \item مورد اول
    \item مورد دوم
    \item مورد سوم

\end{itemize}


% =========================================================

\section{وضعیت تکمیل}

\begin{itemize}

    \item یادگیری مفاهیم پایه
    \item پیاده‌سازی
    \item تست
    \item اتصال به پروژه اصلی

\end{itemize}



















